\chapter{Geometry of numbers}
\section{Completions}
Let \({(K, |\cdot|)}\) a valued field and let \(R\) be the ring of all
Cauchy sequences in \(K\). let \(\mathfrak{m} \subseteq R \) be the ideal of
sequences converging to 0.

\begin{lemma}{}
    \(\mathfrak{m}\) is a maximal ideal in \(R\) 
\end{lemma}
\begin{proof}{}
    Take \({(a_{n})}_n \in  R \sminus \mathfrak{m} \) Since \({(a_{n})}_n
    \not\in \mathfrak{m}\), there exists \(\varepsilon > 0\) and \(N \ge 0\)
    such that \(\forall n \ge N\), \(|a_{n}| \ge \varepsilon\). Now take
    \({(b_{n})}_n\) a sequence such that \(\forall n \ge N\), \(a_{n}b_{n} =
    1\). Then \({(b_{n})}_n \in R\) and hence \({(b_{n})}_n\) in an inverse for
    \({(a_{n})}\) in \(R / \mathfrak{m}\). So \({(R / \mathfrak{m})}\) is a
    field and \(\mathfrak{m}\) is maximal.
\end{proof}

\begin{definition}{}
    The \textbf{completion} of \({(K, |\cdot|)}\) is the field \(\widehat{K} = R
    /\mathfrak{m}\). \(K\) is embedded in \(\widehat{K}\) by sending \(a \in K\)
    to the sequence that is constant equal to \(a\).
\end{definition}

Take \({(a_{n})}_n \in R\). We know \(||a_{m}| - |a_{n}|| \le |a_{m} - a_{n}|\)
so \({(|a_{n}|)}_n \subseteq \mathbb{R} \) is Cauchy. Let \(|{(a_{n})}_n|\) be
its limit. If we take \({(b_{n})}_n \in R\) with \({(a_{n})}_n - {(b_{n})}_n \in
\mathfrak{m}\), then \(||a_{n}| - |b_{n}|| \le |a_{n} - b_{n}|\), so
\({(|a_{n}|)}_n\) and \({(b_{n}|)}_n\) have the same limit. Hence \({(a_{m})}_n
\to |{(a_{n})}_n|\) is an absolute value of \(\widehat{K}\) hat extends the
absolute value of \(K\). If \({(a_{n})}_n \in R\), then \(a_{n} \to
{(a_{n})}_n\) in \(\widehat{K}\).

\begin{proposition}{}
    \({(\widehat{K}, |\cdot |)}\) is complete.
\end{proposition}

\begin{remark}{}
    Let \({(\widehat{K}', |\cdot |')}\) be another complete value field such
    that \(|\cdot |'\) extends the absolute value of \(K\) and \(K\) is dense in
    \(\widehat{K}'\) and \(\widehat{K}'\) is complete. Then there exists a
    unique \(K\)-isomorphism \(\sigma : \widehat{K} \to \widehat{K}'\) such that
    \(|\cdot |' \circ \sigma = |\cdot |\).
\end{remark}

\subsection{Completion of archimedean valued fields}

\begin{theorem}{Ostrowski 2}
    Let \({(K, |\cdot |)}\) be an archimedean complete valued field. There is an
    isomorphism \(\sigma : K \to \mathbb{R}\) or \(K \to \mathbb{C}\) such that 
    \[
      |\cdot | = {(|\cdot|_{\circ} \circ \sigma)}^{s}
    \]
    for some \(\mathbb{R} \ni s \ge 1\) 
\end{theorem}

\subsection{Completion of nonarchimedean valued fields}

\begin{proposition}{}
    Let \({(K, |\cdot|)}\) be a valued field with discrete normalized absolute
    field. Then

\begin{enumerate}[label = \arabic*.]
    \item The valueation on \(\widehat{K}\) is still discrete normalized.
    \item Let \(and {(\widehat{\mathcal{O}})}\) valuation rings of \(K\) and
        \(\widehat{K}\) and \(\mathfrak{p}\) and \(\widehat{\mathfrak{p}}\) the
        maximal ideals. Then
        \[
          \forall n \ge 1 ,\quad \mathcal{O}/\mathfrak{p}^{n} \cong
          \widehat{\mathcal{O}} / \mathfrak{p}^{n}
        \]
        Moreover the diagonal map
        \begin{align*}
            \partial: \widehat{\mathcal{O}} &\longrightarrow
            \lim_{\longleftarrow_{n \ge 1}}  \mathcal{O} / \mathfrak{p}^{n} \\
             a &\longmapsto {(a \mod \mathfrak{p}, a \mod \mathfrak{p}^2,
             \dots)}
        \end{align*}
        is an isomorphism and a homeomrphism (with \(\mathcal{O} /
        \mathfrak{p}^{n}\) endowed with discrete topology).
    \item Let \(\mathcal{R} \subseteq \mathcal{O} \) be a system of
        representatives of the residue field \(k := \mathcal{O} / \mathfrak{p}\)
        containing 0. Let \(\pi\) be a uniformizer of \(K\). Then every element
        \(x \in \widehat{\mathcal{O}}\) can be uniquely written as 
        \[
          x = \sum_{n=0}^{\infty} a_{n} \pi^{n} 
        \]
        with \(a_{0}, a_{1}, \dots \in \mathcal{R}\) 
\end{enumerate}
\end{proposition}

\begin{proof}{}
\begin{enumerate}[label = \arabic*.]
    \item Take \(x \in \widehat{K} \sminus \{0\} \). We can find \(y \in K\)
        with \(|x-y| < |x|\). Then \(|y| = |x|\). Hence \(v {\left(
        \widehat{K}^{\times } \right)} = v {\left( K^{\times } \right)}  \) 
    \item Consider the inclusion \(i : \mathcal{O} \to \widehat{\mathcal{O}}\).
        We have 
        \[
            \mathfrak{p}^{n} = \{x \in \mathcal{O} : v{(x)} \ge n\} = \{x \in
            \widehat{\mathcal{O}} : v{(x)} \ge n\} \cap \mathcal{O} =
            \widehat{\mathfrak{p}}^{n} \cap \mathcal{O}
        \]
        So \(i\) induces an injective morphism \(\widehat{i} : O /
        \mathfrak{p}^{n} \to \widehat{\mathcal{O}} /
        \widehat{\mathfrak{p}}^{n}\). For surjectivity, take \(y \in
        \widehat{\mathcal{O}}\). Since \(K\) is dense in \(\widehat{K}\), we can
        find \({(y_{m})}_m\) with values in \(K\) such that \(y_{m}\to y\).
        \(y + \widehat{\mathfrak{p}}^{n}\) is open. So \(y_{m} \in y +
        \widehat{\mathfrak{p}}^{n}\) for \(m \gg 0\). So \(\widehat{i}{(y_{n})}
        = \widehat{i}{(y)}\) for \(m \gg 0\) and \(\widehat{i}\) is an
        isomorphism.

        Consider the map of the statement \(\partial\). Then \(\mathrm{ker}\,
        \partial = \bigcap_{n \ge 0} \widehat{\mathfrak{p}}^{n} =
        v^{-1}{(\infty)} = 0\). For surjectivity, take \({(\bar{a}_n)}_n\)
        in the limit. Let \(a_{n}\) be a lift of \(\bar{a}_n\) to
        \(\mathcal{O}\) for each \(n\).
        We have \(v{(a_{m} - a_{n})} \ge \min \{m, n\} \), hence \({(a_{n})}_n\)
        is Cauchy. Let \(a\) be its limit in \(\widehat{K}\). Since \(a_{n} \in
        \mathcal{O}\) for all \(n\), \(a \in \widehat{\mathcal{O}}\).
        The inequality \(v{(a_{m} - a_{n})} \ge \min \{m,n\} \) gives \(v{(a -
        a_{n})} \ge n\) for each \(n\), hence \(\partial{(a)} =
        {(\bar{a}_n)}_n\), so \(\partial\) is an isomorphism.

        For the \emph{homeo}, \(\mathfrak{p}^{n}\) form a basis of open
        neighborhoods of \(0\) in \(\widehat{K}\). Then
        \[
          \partial {\left( \mathfrak{p}^{n} \right)} = \mathrm{ker} {\left(
          \lim_{\longleftarrow n} \mathcal{O} / \mathfrak{p}^{n}
  \overset{\text{proj}}{\to} \mathcal{O} / \mathfrak{p}^{n}  \right)} =
  \{{(\bar{0}, \dots, \bar{0}, \bar{a}_{n+1} , \dots)} \in
  \lim_{\longleftarrow n} \mathcal{O} / \mathfrak{p}^{n}\} 
        \]
        which form a basis of open neighborhoods of 0 in the inverse limit,
        hence \(\partial\) is a homeomorphism.

    \item We inductively construct a sequence \({(a_{n})}_n\) with values in
        \(\mathcal{R}\) such that \(\forall m\), \(v{\left( x -
        \sum_{i=0}^{m} a_{i} \pi^{i}  \right)} \ge  m+1\). For \(m=0\), we
        take \(a_{0} \in \mathcal{R}\) such that \(a_{0} \equiv x \mod
        \widehat{\mathfrak{p}}\). Now assume \(a_{m}\) is already constructed.
        Then
        \[
          v {\left( x - \sum_{i=0}^{m} a_{i} \pi^{i}  \right)} \ge m+1 \implies
          x - \sum_{i=0}^{m} a_{i} \pi^{i} = b \pi^{m+1} \quad \text{ with } b
          \in \widehat{\mathcal{O}}
        \]
        Choose \(a_{m+1} \in \mathcal{R}\) such that \(a_{m+1} \equiv b \mod
        \mathfrak{p}\). Then
        \[
          v {\left( x - \sum_{i=0}^{m+1} a_{i} \pi^{i}  \right)} = v {\left(
          \pi^{m+1} {(b - a_{m+1} )} \right)} \ge m+2
        \]
        which finishes the induction.

        Then \(x = \sum_{i=0}^{\infty} a_{i} \pi^{i}\). Take \(b_{i} \in
        \mathcal{R}\) such that \(x = \sum_{i=0}^{\infty} b_{i} \pi^{i} \) and
        assume \({(b_{i})}_i \neq {(a_{i})}_i\).
        Take \(i_{0}\) minimal index with \(b_{i_{0}} \neq a_{i_{0}} \). Then
        \[
          0 = {(a_{i_{0}} - b_{i_{0}} )} \pi^{i_{0}} + \sum_{i \ge b+1}
          {(a_{i} - b_{i} )} \pi^{i}
        \]
        where the first addend has valuation \(i_{0}\) and the second has
        valuation \(\ge i_{0} + 1\). This is a contradiction and thus
        \({(a_{n})}_n\) is unique.
\end{enumerate}

\end{proof}

\begin{definition}{}
    Given \(p\) prime, the completion of \({(\mathbb{Q}, |\cdot |_{v_p} )}\) is
    denoted \(\mathbb{Q}_p\) and is called the \textbf{field of \(p\)-adic
    numbers}. Its valuation ring is denoted \(\mathbb{Z}_p\) and is called the
    \textbf{ring of \(p\)-adic integers}. \(p \mathbb{Z}_p\) is the unique
    maximal ideal of \(\mathbb{Z}_p\) and for all \(n\),
    \[
      \mathbb{Z}_p \cong  \lim_{\longleftarrow n} \mathbb{Z} / p^{n} \mathbb{Z}  \quad \frac{\mathbb{Z}_p}{p^{n}\mathbb{Z}_p} \cong \frac{\mathbb{Z}}{p^{n}\mathbb{Z}} 
    \]
    and an element \(x \in \mathbb{Z}_p\) can be uniquely written as \(x =
    \sum_{i=0}^{\infty} a_{i} p^{i} \), with \(a_{i} \in  \{0, \dots, p-1\} \) 
\end{definition}

\begin{lemmao}{Hensel}
    Let \({(K, |\cdot|)}\) be a complete nonarchimedean valued field with
    valuation ring \(\mathcal{O}\). Take \(f \in  \mathcal{O}[X]\) a polynomial
    and \(a \in \mathcal{O}\) such that \(f'{(a)} \neq 0\) and \(c =
    \frac{|f{(a)}|}{|f'{(a)}|^2} < 1\). Consider the sequence \({(a_{n})}_n\)
    defined as
    \[
      \begin{cases}{}
          a_{0} = a\\
          a_{n+1} = a_{n} - \frac{f{(a_{n})}}{f'{(a_{n})}}
      \end{cases}
    \]
    This sequence is well-defined and converges to a root \(b\) of \(f\) in
    \(\mathcal{O}\) such that \(|b-a| \le c < 1\) 
\end{lemmao}

\begin{eser}{2.2.10}
    
\end{eser}

\subsection{Extensions of complete valued fields}

\begin{definition}{}
    Let \({(K, |\cdot |)}\) be a valued field. Let \(V\) be a finite dimensional
    \(K\)-vector space. A norm \(\|\cdot \|\) on \(V\) is a map \(\|\cdot \| :
    V \to \mathbb{R}_{\ge 0} \) such that \(\forall x \in V\), \(\|x\| = 0 \iff
    x = 0\), it homogeneous and sublinear.

    Another norm \(\|\cdot \|'\) is equivalent to \(\|\cdot \|\) if there exist
    constants \(C_{1}, C_{2} \in \mathbb{R}\) such that \(C_{1}\|\cdot \|' \le
    \|\cdot \| \le C_{2} \|\cdot \|'\) 
\end{definition}

\begin{proposition}{}
    If \(\dim V < \infty\) and \(K\) is complete, then all norms on \(V\) are
    equivalent. In particular, whatever the norm on \(V\), \(V\) is complete.
\end{proposition}
\begin{proof}{}
    Seen a lot of times, just use \(\|\cdot \|_1\).
\end{proof}

\begin{theorem}{}
    Let \({(K, |\cdot |)}\) be a complete valued field. Let \(L\) be a finite
    extension. There exists a unique extension \(|\cdot |_L\) of \(|\cdot |\).
    For all \(x \in L\), 
    \[
    |x|_L = {| N_{L / K} {(x)} |}^{1 / [L:K]}
    \]
    for all \(x \in L\) and \({(L, |\cdot |_L)}\) is complete.
\end{theorem}
\begin{remark}{}
    The archimedean case immediately follows from Ostrowski's theorem.
\end{remark}
\begin{proof}[Proof (assuming \({(K, |\cdot |)}\) discrete
    valued field and \(L / K\)separate)] First let's prove \textbf{existence}.
    Let \(v\) be a valuation on \(K\). The valuation ring \(\mathcal{O}\) of
    \(K\) is a PID. So its integral closure \(B\) in \(L\) is a Dedekind domain.
    We moreover know that \(B\) contains a prime ideal \(\mathfrak{q}\) lying
    above the maximal ideal \(\mathfrak{p}\) of \(\mathcal{O}\). We have the
    \(\mathfrak{q}\)-adic valuation \(v_\mathfrak{q}\) on \(L\), which satisfies
    \( v_\mathfrak{q}|_K = e {(\mathfrak{q}|\mathfrak{p})} v \)
    so \(\frac{v_\mathfrak{q}}{e{(\mathfrak{q}|\mathfrak{p})}}\) is a discrete
    complete valuation on \(L\) that extends \(v\).

    For the \textbf{uniqueness}, let \(|\cdot |_1\) and \(|\cdot|_2\) be two
    absolute values on \(L\) extending \(|\cdot |\). Then \({(L, |\cdot |_1)}\)
    and \({(L,|\cdot |_2)}\) are normed spaces over \({(K, |\cdot |)}\). Since
    \({(K, |\cdot |)}\) is complete and \(L/K\) are finite, \(|\cdot |_1\) and
    \(|\cdot |_2\) are equivalent \emph{norms}, hence equivalent absolute
    values.
    So \(|\cdot|_1 = |\cdot|_2^{s}\) for some \(s\ge 0\). Since \(|\cdot |_1 |_K
    = |\cdot | |_K\), then \(s=1\), hence \(|\cdot |_1 = |\cdot |_2\).

    Let's call \(|\cdot |_L\) the unique extension of \(|\cdot |\) to \(L\). We
    want to prove the equality in the thesis. Let \(M / K\) be the Galois
    closure of \(L / K\), let \(|\cdot |_M\) be the unique extension of \(|\cdot
    |\) to \(M\). For \(\sigma \in \mathrm{Gal}{(M /K)}\), since \(|\cdot |_M\)
    is the \emph{only} extension of \(|\cdot |\) to \(M\), then \(|\cdot |_M
    \circ \sigma = |\cdot |_M\). So, for \(x \in L\),
    \[
        |x|_L = {\left( |x|_L ^{[L:K} \right)}^{1 / [L:K]} = {\left(
        \prod_{\sigma \in \mathrm{Hom}_K {(L, M)}} |\sigma{(x)}|_L \right)} ^{ 1
    / [L:K]} = {|N_{L / K}|}^{1 / [L : K]}
    \]

\end{proof}

\begin{corollary}{}
    Let \({(K, |\cdot |)}\) be a valued field with completion \({(\widehat{K},
    |\cdot |_{\widehat{K}} )}\). Let \(\Omega\) be an algebraic closure of
    \(\widehat{K}\). Let \(L / K\) be a finite extension. Then
\begin{enumerate}[label = \arabic*.]
    \item The absolute value \(|\cdot |_{\widehat{K}}\) has a unique extension
        to \(\Omega\), \(|\cdot|_{\Omega} \) 
    \item The absolute values on \(L\) extending \(|\cdot |\) are the \(|\cdot
        |_\Omega \circ \sigma\), with \(\sigma \in  \mathrm{Hom}_K {(L, \Omega)}\) 
    \item For \(\sigma, \tau \in \mathrm{Hom}_K{(L, \Omega)}\), we have \(|\cdot
        |_\Omega \circ \sigma = |\cdot |_\Omega \circ \tau\) if and only if
        there exists \(\rho \in \mathrm{Gal}{( \Omega / \widehat{K})}\) such
        that \(\tau = \rho \circ \sigma\) 
\end{enumerate}
\end{corollary}
\begin{proof}{} \(\) 
\begin{enumerate}[label = \arabic*.]
    \item For each finite subextension \(M / \widehat{K}\) of \(\Omega /
        \widehat{K}\), the absolute value \(|\cdot |_\widehat{K}\) has a unique
        extension \(|\cdot |_M\) to \(M\). If \(M \subseteq N \) are two finite
        extensions of \(\widehat{K}\) contained in \(\Omega\), then 
        \[
          |\cdot |_N \big|_M = |\cdot |_M
        \]
        This allows to define and absolute value \(|\cdot |_\Omega\) on
        \(\Omega\) extending \(|\cdot |_\widehat{K}\). If \(|\cdot |_{\Omega}
        ^{\prime}\) is another one, then for each \(\Omega / M / \widehat{K}\)
        finite, 
        \[
          |\cdot |^{\prime}_{\Omega} \big|_M = |\cdot |_M = |\cdot |_\Omega
          \big|_M
        \]
        So \(|\cdot |^{\prime}_\Omega = |\cdot |_\Omega\) 
    \item Let \(|\cdot |_L\) be some extension of \(|\cdot |\) on \(L\). Let
        \({\left(\widehat{L}, |\cdot |_\widehat{L}\right)}\) be the completion of \({(L,
        |\cdot |_L)}\). The completion \(\widehat{K}\) of \(K\) identifies with
        the closure of \(K\) in \(\widehat{L}\).
    \[
\begin{tikzcd}
	& {\widehat{L}} & \\
	L & {L\widehat{K}} & {\widehat{K}} \\
	& K
	\arrow[no head, from=2-1, to=1-2]
	\arrow[no head, from=2-1, to=2-2]
	\arrow[no head, from=2-3, to=1-2]
	\arrow[no head, from=2-3, to=2-2]
	\arrow[no head, from=3-2, to=2-1]
	\arrow[no head, from=3-2, to=2-3]
	\arrow[no head, from=2-2, to=1-2]
\end{tikzcd}
    \]
        The know that \(L\widehat{K}\) is complete (finite dimensional over
        \(\widehat{K}\)) and contains \(L\) as a dense subfield, hence
        \(L\widehat{K} = \widehat{L}\).

        So \(\widehat{L} / \widehat{K}\) is finite and we can find \(\sigma \in
        \mathrm{Hom}_\widehat{K} {(\widehat{L}, \Omega)}\) such that \(|\cdot
        |_\widehat{L}\) and \(|\cdot |_\Omega \circ \sigma\) are absolute values
        on \(\widehat{L}\) extending \(|\cdot |_\widehat{K}\). So
        \[
        |\cdot |_{\widehat{L}} = |\cdot |_\Omega \circ \sigma \quad \text{ and }
          \quad |\cdot |_L = |\cdot |_\Omega \circ \sigma|_L
        \]
    \item Assume \(\sigma = \rho \circ \tau\), with \(\rho \in
        \mathrm{Aut}{(\Omega / \widehat{K})}\). Then \(|\cdot |_\Omega \circ
        \sigma = |\cdot |_\Omega \circ \rho \circ \tau\). But \(|\cdot |_\Omega
        \circ \rho\) is an absolute value on \(\Omega\). So, by 1), \(|\cdot
        |_\Omega \circ \rho = |\cdot |_\Omega\). So \(|\cdot |_\Omega \circ
        \sigma = |\cdot |_\Omega \circ \tau\).

        Conversely, assume \(|\cdot |_\Omega \circ \sigma = |\cdot |_\Omega
        \circ \tau\). Then set \(\rho_0 := \tau \circ \sigma^{-1} : \sigma {(L)}
        \to \tau {(L)}\). We claim that \(\rho_{0}\) can be extended to some
        \(\widehat{\rho_{0}} : \sigma{(L)} \widehat{K} \to \tau{(L)}
        \widehat{K}\). To do so, take \(x \in \sigma{(L)}\widehat{K} \to \tau
        {(L)} \widehat{K}\) continuously.
        To do so, take \(x \in \sigma{(L)}\widehat{K}\) and choose
        \({(x_{n})}_n\) a sequence in \(\sigma{(L)}\) converging to \(x\). For
        \(m, n \ge 1\), 
        \begin{align*}
            |\rho_{0}{(x_{m})} - \rho_{0}{(x_{n})}|_{\Omega} &= |\tau \circ
            \sigma^{-1}{(x_{m})} - \tau \circ \sigma^{-1}{(x_{n})}|_\Omega =
            |\tau \circ \sigma^{-1}{(x_{m} - x_{n})}|_{W} = \\
             &= |\sigma \circ \sigma^{-1} {(x_{m} - x_{n})} |_\Omega = |x_{m} -
             x_{n}|_\Omega
        \end{align*}
        Hence \({(\rho_{0}{(x_{n})})}_n\) is Cauchy, hence it has a limit
        \(\rho_{0}{(x)}\) in \(\sigma{(L)}{(\widehat{K})}\) (which is complete).

        One then checks that \(\widehat{\rho_{0}}{(x)}\) only depends on \(x\)
        and that \(\widehat{\rho_{0}}\) is a \(\widehat{K}\)-isomorphism
        (exercise). We end by extending \(\widehat{\rho_{0}}\) to some \(\rho
        \in \mathrm{Aut}{(\Omega / \widehat{K})}\). Then we can conclude that
        \[
          \sigma \circ \rho = \tau
        \]
    
\end{enumerate}

\end{proof}

\section{Ramification Theory}
\subsection{Complete discrete valued fields}

Here we fix \({(K, |\cdot |_K)}\) a complete discretely valued field. We denote
by \(v_K\) the associated valuation. Let \(L\) be a finite extension of \(K\).
We denote by \(|\cdot | = |\cdot |_L\) the extension of \(|\cdot |_K\) to \(L\)
and \(v = v_L\) the extension of \(v_K\) to \(L\).

We assume \(L / K\) is \emph{separable} (not always needed, but leads to easier
proofs).

\begin{lemma}{}
    The valuation ring \(\mathcal{O}_L\) of \(L\) is the integral closure of the
    valuation ring \(\mathcal{O}_K\) of \(K\) in \(L\).
\end{lemma}

\begin{proof}{}
    Since \(\mathcal{O}_L\) is integrally closed, it contains the integral
    closure of \(\mathcal{O}_K\) in \(L\). Take \(x \in \mathcal{O}_L\). Since
    \(|\cdot |_L = {\left( |\cdot |_K \circ N_{L / K} \right)}^{1 / [L : K]} \).
    We have
    \[1 \ge  |N_{L / K} {(x)}|_K = \left| \prod_{\sigma \in \mathrm{Hom}_{K}
    {(L, K^{\text{sep}})}} \sigma{(x)} \right|_K \] 

    Let \(\Omega\) be an algebraic closure of \(K\). Let \(|\cdot |_\Omega\) be
    the extension of \(|\cdot |_K\) to \(\Omega\). By uniqueness, for \(\sigma
    \in \mathrm{Hom}_K{(L, \Omega)}\), we have that \(|\sigma{(x)}|_\Omega =
    |x|_{\Omega} \). So \(N_{L/K}{(x)}|_K = |x|_L^{[L : K]} \), and so
    \(|\sigma{(x)}|_L \le 1\) for all \(\sigma \in \mathrm{Hom}_K {(L,
    \Omega)}\).

    Using the relations of roots and coefficients, we deduce that the minimal
    polynomial \(\mu_x\) of \(x\) is in \(\mathcal{O}_K[X]\), so \(x\) is
    integral over \(\mathcal{O}_K\).
\end{proof}

In the following, we will use the following notations:
\begin{itemize}[label = --]
    \item \(\mathcal{O}_K, \mathcal{O}_L\) valuation rings
    \item \(\mathfrak{p}_K, \mathfrak{p}_L\) maximal ideals
    \item \(\kappa_K, \kappa_L\) the residue fields
    \item \(e = e{(L/K)} = e {(\mathfrak{p}_L | \mathfrak{p}_K)}\) the
        ramification index: \(\mathfrak{p}_K \mathcal{O}_L = \mathfrak{p}_L^{e}\) 
    \item \(f = f{(L / K)} = f{(\mathfrak{p}_L | \mathfrak{p}_K)}\) the residual
        degree \(f = [\kappa_L : \kappa_K]\) 
\end{itemize}
\begin{remark}{}
    \(e = \left[ v_L{(L^{\times })} : v_K{(\kappa^{\times })} \right] \) 
\end{remark}
And recall that \(ef = [L : K]\), that \(L / K\) unramified means that
\(\kappa_L / \kappa_K\) is separable and \(e=1\), and that \(L / K\) is totally
ramified if \(f=1\) 

\begin{lemma}{}
    Take \(a \in \mathcal{O}_L\). Assume that \(f{(a)} = 0\) for some monic \(f
    \in \mathcal{O}_K[X]\) and that \(\bar{a} \in \kappa_L\) is a simple root of
    \(\bar{f} \in \kappa_K[X]\). Then \(\kappa{(a)} / \kappa\) is unramified,
    i.e.
    \[
      \kappa_{K{(a)}} = \kappa_K {(\bar{a})} \quad \text{ and } \quad \left[
      \kappa_{K{(a)}} : \kappa_K \right] = \left[ K{(a)} : K \right] 
    \]
\end{lemma}
\begin{proof}{}
    Let \(\bar{g} \in \kappa_K[X]\) be the minimal polynomial of \(\bar{a}\).
    Let \(g \in \mathcal{O}_K[X]\) be a monic lifting. Since \(\bar{a}\) is a
    simple root of \(\bar{f}\), it is also a simple root of \(\bar{g}\). It is
    separable over \(\kappa_K\), so \(\kappa_K{(\bar{a})} / \kappa_K\) is
    separable, and by Hensel's lemma we can find \(b \in \mathcal{O}_L\) such
    that \(g{(b)} = 0\) and \(\bar{b} = \bar{a}\).
    \[
      \deg g = \deg \bar{g} \le  \left[ \kappa_{K{(b)}} : \kappa \right] \le  \left[
      K{(b)} : K \right] \le  \deg g
    \]
    So \([K{(b)} : K] = \left[ \kappa_{K{(b)}} : \kappa_K \right] \). So
    \(K{(b)} / K\) is unramified and \(\kappa_{K{(b)}} = \kappa{(\bar{a})}\).
    Applying Hensel's lemma to \(f\) over \(K{(b)}\), we can find a root \(c \in
    \mathcal{O}_{K{(b)}} \) of \(f\) lifting \(\bar{a}\). Such a root is
    necessarily unique. So we get that \(c=a\), hence \(a \in K{(b)}\).

    Finally
    \[
      \deg \bar{g} \overset{\bar{a} \in \kappa_{K{(a)}} }{\le }  \left[
      \kappa_{K{(a)}} : \kappa_K \right] \le \left[ K{(a)} : K \right] \le
      \left[ K{(b)} : K \right] = \deg g
    \]
    Hence \(K{(a)} = K{(b)}\) 
\end{proof}

